\chapter{Приложения}

\section*{Краткий испано-русский словарь путешественника}
\addcontentsline{toc}{section}{Словарь}

\begin{multicols}{2}\small\raggedright
\textbf{ayuntamiento} --- ратуша \\
\textbf{barranco} --- ущелье \\
\textbf{calima} --- пыльный ветер из~Сахары \\
\textbf{caldera} --- вулканическая котловина \\
\textbf{casa cueva} --- пещерный дом \\
\textbf{degollada} --- седловина \\
\textbf{ermita} --- часовня \\
\textbf{gofio} --- обжаренная мука \\
\textbf{guanche} --- доиспанский житель \\
\textbf{guagua} --- автобус \\
\textbf{laurisilva} --- туманный лавровый лес \\
\textbf{malpaís} --- «дурная земля», лавовое поле \\
\textbf{mirador} --- смотровая площадка \\
\textbf{mojo} --- соус на~оливковом масле \\
\textbf{pinar} --- сосновый бор \\
\textbf{playa} --- пляж \\
\textbf{risco} --- отвесная скала \\
\textbf{roque} --- скальная игла \\
\textbf{sendero} --- тропа \\
\textbf{tagoror} --- место гуанчского совета \\
\end{multicols}

\section*{Источники и~полезные ссылки}
\addcontentsline{toc}{section}{Источники}

\begin{itemize}
\item Картографическая основа: IGN Espana 1:25\,000, листы 1101-III,
      1109-I, 1110-IV.
\item Расписание автобусов: приложение \textit{Global Transporte},
      \url{https://guaguasglobal.com}.
\item Метеопрогноз: AEMET Canarias, \url{https://www.aemet.es}.
\item Маркированные тропы: портал \textit{Gran Canaria Camina},
      \url{https://grancanariacamina.com}.
\item Охраняемые территории: сеть \textit{Red Canaria de~Espacios
      Naturales Protegidos}.
\end{itemize}

\vspace{1em}
\noindent{\color{baedmute}\small\itshape
Сведения проверены по~состоянию на~сезон 2024/25; перед поездкой
рекомендуется уточнить актуальность расписаний и~цен.}
